The results identify two service objectives that require coordination:
making original computation timely and keeping its supported progress
available. First-delivery credit captures the former event; whole-view
support captures the latter state. Their ranking can differ, as the
continuity experiment shows. A service contract therefore needs both a
deadline and a rule for what remains available after feedback changes.

The native comparison isolates a second choice within the same safety
contract: which issued work remains available for reuse. Ready-only and
Fullcall share original selection and view-support checks, but only the
latter preserves a call while its input is being resolved. Input-bound
versions can supply the same eligibility evidence. This gives builders
freedom in representing dependencies while retaining the separate
checks at execution, selection, and application.

\paragraph{What an implementation must preserve}
Three obligations survive any encoding. First, input eligibility:
a representation supports safe in-flight retention only if it can witness
the complete input of the call at the installed parent; a coarse birth
version cannot. Second, original selection: retaining executed work never
endorses its prediction, and the original Correct/Update and the durable
selection boundary stay unchanged. Third, view support: protecting
computation does not protect delivery, so the Consumer validates the
whole view at each actual application. The contract is informational, not
representational: the measured version-DAG policy matches Fullcall's
eligible handovers.

\paragraph{Operating scope}
The measured gains separate a correctness contract from the scheduling
conditions that make reuse profitable: witnesses govern the contract,
while available execution overlap and Consumer timing govern the
profit. Local agreement reduces revisions by waiting for more
evidence; independent canonical chains can consume the slots needed for
conditional execution~\cite{welsh2001seda}. The selected-fact
argument establishes allowed transitions under the execution contract,
while controlled fault injections exercise its lifecycle boundaries.

Adapting this protocol requires a closed, serializably comparable call
boundary: all state read by the original inference and feedback update,
its complete result and selected transition, and the published
dependencies of the Consumer's current view. These interfaces make
eligibility and continued support checkable at the point where retained
work is used. Numerical cross-call state, such as a reused KV cache,
would need an approximate or prefix-matched witness before the same
checks apply.
